% =========================================================================
% Figura: Interpretación geométrica del indicador de Espaciamiento (Spacing)
% Problema de minimización bidimensional
% =========================================================================

\begin{figure}[htbp]
\centering
\resizebox{0.58\linewidth}{!}{%
\begin{tikzpicture}[
    font=\sffamily,
    >=Stealth,
    x=1cm, y=1cm
]

% -------------------------------------------------------------------------
% 1. EJES COORDENADOS (Minimización)
% -------------------------------------------------------------------------
\draw[line width=1.1pt, draw=darkText, ->, -{Stealth[length=6pt, width=4pt]}]
    (0, 0) -- (6.8, 0) node[below=5pt, midway, font=\fontsize{9.5}{12}\selectfont, text=darkText] {\textbf{Objetivo $f_1$} (minimizar)};

\draw[line width=1.1pt, draw=darkText, ->, -{Stealth[length=6pt, width=4pt]}]
    (0, 0) -- (0, 5.7);
\node[font=\fontsize{9.5}{12}\selectfont, text=darkText, rotate=90]
    at (-0.65, 2.85) {\textbf{Objetivo $f_2$} (minimizar)};

% -------------------------------------------------------------------------
% 2. TRAYECTORIA DEL FRENTE APROXIMADO (Línea suave discontinua)
% -------------------------------------------------------------------------
\draw[draw=grayText!45, dashed, line width=0.8pt]
    (0.9, 4.9) to[out=-70, in=150] (1.90, 3.20) to[out=-30, in=165] (4.00, 1.55) to[out=-15, in=175] (6.0, 0.9);

% Coordenadas de las soluciones obtenidas A
\coordinate (a1) at (1.20, 4.40);
\coordinate (a2) at (1.90, 3.20);
\coordinate (a3) at (2.85, 2.25);
\coordinate (a4) at (4.00, 1.55);
\coordinate (a5) at (5.40, 1.05);

% -------------------------------------------------------------------------
% 3. DISTANCIAS ENTRE SOLUCIONES CONSECUTIVAS d_i (Flechas dobles bermellón)
% -------------------------------------------------------------------------
\draw[line width=1.0pt, draw=figVermilion, <->, >={Stealth[length=4pt, width=2.8pt]}, shorten <=0.14cm, shorten >=0.14cm]
    (a1) -- (a2) node[midway, above right=1pt, font=\fontsize{7.5}{9}\selectfont\bfseries, text=figVermilion!90!black] {$d_1$};

\draw[line width=1.0pt, draw=figVermilion, <->, >={Stealth[length=4pt, width=2.8pt]}, shorten <=0.14cm, shorten >=0.14cm]
    (a2) -- (a3) node[midway, above right=1pt, font=\fontsize{7.5}{9}\selectfont\bfseries, text=figVermilion!90!black] {$d_2$};

\draw[line width=1.0pt, draw=figVermilion, <->, >={Stealth[length=4pt, width=2.8pt]}, shorten <=0.14cm, shorten >=0.14cm]
    (a3) -- (a4) node[midway, above right=1pt, font=\fontsize{7.5}{9}\selectfont\bfseries, text=figVermilion!90!black] {$d_3$};

\draw[line width=1.0pt, draw=figVermilion, <->, >={Stealth[length=4pt, width=2.8pt]}, shorten <=0.14cm, shorten >=0.14cm]
    (a4) -- (a5) node[midway, above right=1pt, font=\fontsize{7.5}{9}\selectfont\bfseries, text=figVermilion!90!black] {$d_4$};

% Dibujar soluciones A (círculos azules)
\newcommand{\drawPointSpacing}[1]{%
    \filldraw[draw=figBlue!90!black, fill=figBlue, line width=0.9pt] (#1) circle (0.12cm);
}
\drawPointSpacing{a1}
\drawPointSpacing{a2}
\drawPointSpacing{a3}
\drawPointSpacing{a4}
\drawPointSpacing{a5}

% -------------------------------------------------------------------------
% 4. LEYENDA (Esquina superior derecha libre)
% -------------------------------------------------------------------------
\begin{scope}[shift={(2.90, 3.75)}]
    \draw[draw=grayText!40, fill=white, rounded corners=3pt, line width=0.6pt]
        (0, 0) rectangle (3.75, 1.30);

    % Soluciones obtenidas A
    \filldraw[draw=figBlue!90!black, fill=figBlue, line width=0.6pt] (0.20, 1.00) circle (0.08cm);
    \node[anchor=west, font=\fontsize{6.5}{7.8}\selectfont, text=darkText]
        at (0.36, 1.00) {Soluciones obtenidas ($A$)};

    % Distancia relativa d_i
    \draw[line width=1.0pt, draw=figVermilion, <->, >={Stealth[length=3.5pt, width=2.4pt]}]
        (0.08, 0.65) -- (0.32, 0.65);
    \node[anchor=west, font=\fontsize{6.5}{7.8}\selectfont, text=darkText]
        at (0.36, 0.65) {Distancia relativa ($d_i$)};

    % Frente aproximado
    \draw[draw=grayText!60, dashed, line width=0.7pt] (0.08, 0.30) -- (0.32, 0.30);
    \node[anchor=west, font=\fontsize{6.5}{7.8}\selectfont, text=darkText]
        at (0.36, 0.30) {Frente aproximado};
\end{scope}

\end{tikzpicture}%
}
\caption{Interpretación geométrica del indicador de espaciamiento (\textit{Spacing}) en un espacio de dos objetivos a minimizar. Mide la desviación estándar de las distancias euclídeas $d_i$ entre soluciones consecutivas a lo largo del frente aproximado. Valores próximos a cero indican una distribución uniforme sin acumulaciones ni discontinuidades.}
\label{fig:indicador_espaciamiento}
\end{figure}
